\section{从论文架构到现实系统：BAGEL 与 \texorpdfstring{$\pi_{0.7}$}{pi0.7}}

MoT 的意义不只是一篇 loss-curve 论文，而是提供了一个可组合的 design pattern：让不同 modality 在同一序列中交流，但为生成、理解或控制分配专用参数。课堂选择 BAGEL 与 $\pi_{0.7}$，分别展示这个模式怎样进入统一图文预训练和 embodied control。两者不是同一任务的 benchmark 对手，而是同一架构思想在不同 capability frontier 上的实例。

\subsection{BAGEL：理解共享，生成专用}

前一章建立了 modality-aware routing，本节看它在统一图文模型中的实际组合方式。BAGEL 不是为了证明所有模块都能合并，而是展示怎样让 understanding 保持共享、generation 保留专用路径。

\lecturefigure{slide-045.jpg}{BAGEL 将 modality-aware 参数用于统一理解与生成}{官方 deck 第 45 页；BAGEL arXiv 2505.14683v3；课堂 00:31:39--00:33:33}

\begin{knowledgebox}{读图：BAGEL 不是简单的“一个 encoder 加一个 diffusion head”}
slide 左侧列出 understanding、generation 与 editing 等任务，右侧论文图强调 unified pretraining。课堂解释的关键是参数路径：image generation 使用专用参数，而理解侧的 multimodal language model 仍可共同处理 text 与 image input。这样既保留 language reasoning，又避免要求 generative latent 与 semantic input embedding 完全相同。
\end{knowledgebox}

BAGEL 还展示了一个重要顺序：模型可以先在 text space 中规划，再生成 image。可把过程抽象成
$$
p(v\mid t)=\sum_r p_\theta(r\mid t)\,p_\theta(v\mid t,r),
$$
其中 $r$ 是可见或隐式 reasoning/planning representation，$v$ 是视觉输出。这个分解不是说所有图像都必须先写一段 chain-of-thought，而是说明 language abstraction 可以作为控制与组合接口。

\begin{importantbox}{BAGEL 提供的设计启示}
Understanding 与 generation 可以共享高层语义和 sequence context，同时使用不同 image representation 与参数路径。统一系统的价值来自任务间可组合，而不是追求每一层、每一 loss、每一 token 都同构。
\end{importantbox}

\teachervoice{00:32:51--00:33:33，老师再次把当前状态描述为“尚未找到有效方式完全统一 image generation 与 understanding”。BAGEL 是一种实用分工，不是这个 open problem 已经被理论解决。}

\subsection{\texorpdfstring{$\pi_{0.7}$}{pi0.7}：多模态 context 不只是语言指令}

从数字内容生成继续向前，下一步是让 multimodal context 影响真实动作。本节以 $\pi_{0.7}$ 说明 steering signal 可以包含 subgoal image 与 episode metadata，并讨论为何 embodiment 需要不同评价闭环。

\lecturefigure{slide-046.jpg}{$\pi_{0.7}$ 把 modality-aware foundation model 推向机器人控制}{官方 deck 第 46 页；$\pi_{0.7}$ arXiv 2604.15483v2；课堂 00:34:00--00:35:08}

\begin{knowledgebox}{读图：从生成内容转向生成动作}
图中展示多种 robot、task 与 context 条件。$\pi_{0.7}$ 的 steering 不只依赖一句 language command，还可以利用 subgoal image、episode quality 或策略信息。与图像生成相比，机器人输出的评价更依赖 task success、速度、碰撞和跨 embodiment generalization；一个动作看起来合理，不等于闭环执行成功。
\end{knowledgebox}

若 policy 以 observation $o_t$、language instruction $g$ 与 multimodal steering context $c$ 为条件，则
$$
a_t\sim\pi_\theta(a_t\mid o_{\le t},g,c).
$$
$c$ 可以编码“怎样做”而不只是“做什么”。这说明非文本 modality 的价值不仅是提供更多感知，也可以成为 strategy specification。

\begin{warningbox}{Embodiment evidence boundary}
$\pi_{0.7}$ 的机器人结果支持多任务、跨环境与 steerability，但不能从 slide 推出通用 physical intelligence。真实部署还要求 control frequency、sensor calibration、failure recovery、安全约束和 hardware-specific latency。课堂把它列为 MoT-style specialization 的实例，而非宣告机器人问题已经解决。
\end{warningbox}

\subsection{本章小结}

BAGEL 与 $\pi_{0.7}$ 表明 modality-aware architecture 可以向两个方向扩展：一是统一内容理解、编辑和生成，二是把 multimodal context 变成动作策略条件。共同规律是\textbf{允许高层交互、保留低层专用性}。这也把问题带回讲座的第二条主线：理解、生成与行动之间究竟是否存在稳定的 positive transfer？

